% Source: Weinberg GR, physical PDF pp. 175--179
%         (printed pp. 151--155).
% Coverage: Section 7.1, Derivation of the Field Equations; equations
%           (7.1.1)--(7.1.16).
% Figures/tables/footnotes: reference marker 1; no figures or tables.
% Status: source-reviewed and compile-clean.
% Uncertainties: none.
% MODERNIZATION: The source curvature convention is the negative of the
% binding convention at fixed slot order. The weak-field signs, trace, and
% cosmological-constant equation have therefore been rederived for
% G_{\mu\nu}+\Lambda g_{\mu\nu}=8\pi G T_{\mu\nu}.

\subsection{Derivation of the Field Equations}\label{sec:7.1}

The field equations for gravitation are inevitably going to be more
complicated than those for electromagnetism. Maxwell's equations are linear
because the electromagnetic field does not itself carry charge, whereas
gravitational fields do carry energy and momentum (see
Section~\ref{sec:5.3}) and must therefore contribute to their own source.
That is, the gravitational field equations will have to be nonlinear partial
differential equations, the nonlinearity representing the effect of
gravitation on itself.

In dealing with these nonlinear effects we are guided once again by the
Principle of Equivalence. At any point \(X\) in an arbitrarily strong
gravitational field, we can define a locally inertial coordinate system such
that
\begin{equation}
  g_{\mu\nu}(X)=\eta_{\mu\nu}
  \tag{7.1.1}\label{eq:7.1.1}
\end{equation}
and
\begin{equation}
  \left.\partial_\rho g_{\mu\nu}(x)\right|_{x=X}=0.
  \tag{7.1.2}\label{eq:7.1.2}
\end{equation}
Hence for \(x\) near \(X\), the metric tensor \(g_{\mu\nu}\) can differ from
\(\eta_{\mu\nu}\) only by terms quadratic in \(x-X\). In this coordinate
system the gravitational field is weak near \(X\), and we can hope to
describe the field by linear partial differential equations. And once we
know what these weak-field equations are, we can find the general field
equations by reversing the coordinate transformation that made the field
weak.

Unfortunately, we have very little empirical information about the weak-field
equations. This is not for any fundamental reason, but rather because
gravitational radiation is so weakly generated and absorbed by matter that it
has not yet certainly been detected. However, although forgivable, our
ignorance does prevent us from proceeding as directly as we did in previous
chapters, and some guesswork will be unavoidable.

First let us recall that in a weak static field produced by a nonrelativistic
mass density \(\rho\), the time-time component of the metric tensor is
approximately given by
\[
  g_{00}\simeq-(1+2\phi).
\]
(See Eq.~\eqref{eq:3.4.5}.) Here \(\phi\) is the Newtonian potential,
determined by Poisson's equation
\[
  \boldsymbol{\nabla}^{2}\phi=4\pi G\rho,
\]
where \(G\) is Newton's constant, equal to
\(6.670\times10^{-8}\) in c.g.s.\ units. Furthermore, the energy density
\(T_{00}\) for nonrelativistic matter is just equal to its mass density,
\[
  T_{00}\simeq\rho.
\]
Combining the above, we have then
\begin{equation}
  \boldsymbol{\nabla}^{2}g_{00}=-8\pi G T_{00}.
  \tag{7.1.3}\label{eq:7.1.3}
\end{equation}
This field equation is only supposed to hold for weak static fields generated
by nonrelativistic matter, and is not even Lorentz invariant as it stands.
However, Eq.~\eqref{eq:7.1.3} leads us to guess that the weak-field equations
for a general distribution \(T_{\mu\nu}\) of energy and momentum take the
form
\begin{equation}
  \mathcal G_{\mu\nu}=8\pi G T_{\mu\nu},
  \tag{7.1.4}\label{eq:7.1.4}
\end{equation}
where \(\mathcal G_{\mu\nu}\) is a linear combination of the metric and its
first and second derivatives. It follows then from the Principle of
Equivalence that the equations which govern gravitational fields of arbitrary
strength must take the form
\begin{equation}
  \mathcal G_{\mu\nu}=8\pi G T_{\mu\nu},
  \tag{7.1.5}\label{eq:7.1.5}
\end{equation}
where \(\mathcal G_{\mu\nu}\) is a tensor which reduces to the weak-field
\(\mathcal G_{\mu\nu}\) of Eq.~\eqref{eq:7.1.4}.

In general, there will be a variety of tensors \(\mathcal G_{\mu\nu}\) that
can be formed from the metric tensor and its derivatives, and that reduce in
the weak-field limit to a given \(\mathcal G_{\mu\nu}\). Let us imagine
\(\mathcal G_{\mu\nu}\) to be expanded in a sum of products of derivatives of
the metric, and classify each term according to the total number \(N\) of
derivatives of metric components. (For example, a term with \(N=3\) could be
linear in third derivatives of the metric, or a product of a first derivative
with a second derivative, or a product of three first derivatives.) The whole
of \(\mathcal G_{\mu\nu}\) must have the dimensions of a second derivative, so
each term of type \(N\ne2\) appears multiplied by a constant having the
dimensions of length to the power \(N-2\); such terms will become negligible
for gravitational fields of sufficiently large or small spacetime scale if
\(N>2\) or \(N<2\), respectively. In order to remove the ambiguity in
\(\mathcal G_{\mu\nu}\), we shall assume that \emph{the gravitational field
equations are uniform in scale}, so that only terms with \(N=2\) are allowed.

Let us review what we know about the left-hand side of the field equation
\eqref{eq:7.1.5}:

\begin{description}
  \item[(A)] By definition, \(\mathcal G_{\mu\nu}\) is a tensor.

  \item[(B)] By assumption, \(\mathcal G_{\mu\nu}\) consists only of terms
  with \(N=2\) derivatives of the metric; that is,
  \(\mathcal G_{\mu\nu}\) contains only terms that are either linear in the
  second derivatives or quadratic in the first derivatives of the metric.

  \item[(C)] Since \(T_{\mu\nu}\) is symmetric, so is
  \(\mathcal G_{\mu\nu}\).

  \item[(D)] Since \(T_{\mu\nu}\) is conserved (in the sense of covariant
  differentiation), so is \(\mathcal G_{\mu\nu}\):
\end{description}
\begin{equation}
  \nabla_\mu\mathcal G^{\mu\nu}=0.
  \tag{7.1.6}\label{eq:7.1.6}
\end{equation}
\begin{description}
  \item[(E)] For a weak stationary field produced by nonrelativistic matter,
  the \(00\) component of Eq.~\eqref{eq:7.1.5} must reduce to
  Eq.~\eqref{eq:7.1.3}, so in this limit
\end{description}
\begin{equation}
  \mathcal G_{00}\simeq-\boldsymbol{\nabla}^{2}g_{00}.
  \tag{7.1.7}\label{eq:7.1.7}
\end{equation}
These properties are all we will need to find \(\mathcal G_{\mu\nu}\).

We saw in Section~\ref{sec:6.2} that the most general way of constructing a
field satisfying (A) and (B) is by contraction of the curvature tensor
\(R^\rho{}_{\sigma\mu\nu}\). The antisymmetry property of
\(R_{\rho\sigma\mu\nu}\), discussed in Section~\ref{sec:6.6}, shows that
there are only two tensors that can be formed by contracting
\(R_{\rho\sigma\mu\nu}\); that is, the Ricci tensor
\(R_{\mu\nu}=R^\rho{}_{\mu\rho\nu}\), and the curvature scalar
\(R=R^\mu{}_\mu\). Hence (A) and (B) require
\(\mathcal G_{\mu\nu}\) to take the form
\begin{equation}
  \mathcal G_{\mu\nu}
  =C_1R_{\mu\nu}+C_2g_{\mu\nu}R,
  \tag{7.1.8}\label{eq:7.1.8}
\end{equation}
where \(C_1\) and \(C_2\) are constants. This is automatically symmetric
(see Eq.~\eqref{eq:6.6.7}), so (C) tells us nothing new. Using the Bianchi
identity \eqref{eq:6.8.3} gives the divergence of
\(\mathcal G_{\mu\nu}\) as
\[
  \nabla_\mu\mathcal G^{\mu\nu}
  =
  \left(\frac{C_1}{2}+C_2\right)\nabla^\nu R,
\]
so (D) allows two possibilities: either \(C_2=-C_1/2\), or
\[
  \mathcal G^\mu{}_\mu=(C_1+4C_2)R=\text{constant}.
\]
We can reject the second possibility, because Eq.~\eqref{eq:7.1.5} then says
that if \(\nabla_\nu R=0\), then so must
\(\nabla_\nu T^\mu{}_\mu\), and this is not the case in the presence of
inhomogeneous nonrelativistic matter. We conclude then that
\(C_2=-C_1/2\), so Eq.~\eqref{eq:7.1.8} becomes
\begin{equation}
  \mathcal G_{\mu\nu}
  =
  C_1\left(R_{\mu\nu}-\frac12g_{\mu\nu}R\right).
  \tag{7.1.9}\label{eq:7.1.9}
\end{equation}

Finally, we use the property (E) to fix the constant \(C_1\). A
nonrelativistic system always has \(\lvert T_{ij}\rvert\ll\lvert
T_{00}\rvert\), so we are concerned here with a case where
\(\lvert\mathcal G_{ij}\rvert\ll\lvert\mathcal G_{00}\rvert\). Using
Eq.~\eqref{eq:7.1.9},
\[
  R_{ij}\simeq\frac12\delta_{ij}R.
\]
Furthermore, we deal here with a weak field, so
\(g_{\mu\nu}\simeq\eta_{\mu\nu}\). The curvature scalar is therefore given
by
\[
  R\simeq R_{ii}-R_{00}\simeq\frac32R-R_{00},
\]
or
\begin{equation}
  R\simeq2R_{00}.
  \tag{7.1.10}\label{eq:7.1.10}
\end{equation}
Using Eqs.~\eqref{eq:7.1.10} and \eqref{eq:7.1.1} in
Eq.~\eqref{eq:7.1.9}, we find
\begin{equation}
  \mathcal G_{00}\simeq2C_1R_{00}.
  \tag{7.1.11}\label{eq:7.1.11}
\end{equation}
To calculate \(R_{00}\) for a weak field we may use the linear part of
\(R_{\rho\sigma\mu\nu}\), given by Eq.~\eqref{eq:6.6.2}, as
\[
  R_{\rho\sigma\mu\nu}
  \simeq
  \frac12\left(
    \partial_\sigma\partial_\mu g_{\rho\nu}
    +\partial_\rho\partial_\nu g_{\sigma\mu}
    -\partial_\sigma\partial_\nu g_{\rho\mu}
    -\partial_\rho\partial_\mu g_{\sigma\nu}
  \right).
\]
When the field is static all time derivatives vanish, and the components we
need become
\[
  R_{0000}=0,
  \qquad
  R_{i0j0}\simeq-\frac12\partial_i\partial_jg_{00}.
\]
Hence Eq.~\eqref{eq:7.1.11} gives
\[
  \mathcal G_{00}
  \simeq2C_1(R_{i0\ii0}-R_{0000})
  \simeq-C_1\boldsymbol{\nabla}^{2}g_{00},
\]
and comparing this with Eq.~\eqref{eq:7.1.7}, we find that (E) is satisfied
if and only if \(C_1=1\). Setting \(C_1=1\) in Eq.~\eqref{eq:7.1.9}
completes our calculation of \(\mathcal G_{\mu\nu}\):
\begin{equation}
  G_{\mu\nu}
  \equiv
  R_{\mu\nu}-\frac12g_{\mu\nu}R.
  \tag{7.1.12}\label{eq:7.1.12}
\end{equation}
With Eq.~\eqref{eq:7.1.5}, this gives the \emph{Einstein field equations}
\begin{equation}
  R_{\mu\nu}-\frac12g_{\mu\nu}R
  =
  8\pi G T_{\mu\nu}.
  \tag{7.1.13}\label{eq:7.1.13}
\end{equation}
An alternative form is sometimes useful. Contracting
Eq.~\eqref{eq:7.1.13} with \(g^{\mu\nu}\) gives
\[
  R-2R=8\pi G T^\mu{}_\mu,
\]
or
\begin{equation}
  R=-8\pi G T^\mu{}_\mu,
  \tag{7.1.14}\label{eq:7.1.14}
\end{equation}
and using this in Eq.~\eqref{eq:7.1.13}, we have
\begin{equation}
  R_{\mu\nu}
  =
  8\pi G\left(
    T_{\mu\nu}-\frac12g_{\mu\nu}T^\lambda{}_\lambda
  \right).
  \tag{7.1.15}\label{eq:7.1.15}
\end{equation}
Of course we can also go from Eq.~\eqref{eq:7.1.15} back to
Eqs.~\eqref{eq:7.1.14} and \eqref{eq:7.1.13}, so
Eqs.~\eqref{eq:7.1.13} and \eqref{eq:7.1.15} should be regarded as entirely
equivalent forms of the Einstein field equations.

In a vacuum \(T_{\mu\nu}\) vanishes, so from Eq.~\eqref{eq:7.1.15} we see
that the Einstein field equations in \emph{empty space} are just
\begin{equation}
  R_{\mu\nu}=0.
  \tag{7.1.16}\label{eq:7.1.16}
\end{equation}
In a spacetime of two or three dimensions this would imply the vanishing of
the full curvature tensor \(R_{\rho\sigma\mu\nu}\), and the consequent
absence of a gravitational field. (See Section~\ref{sec:6.4}.) It is only in
four or more dimensions that true gravitational fields can exist in empty
space.

We might be willing to relax assumption (B), and allow
\(\mathcal G_{\mu\nu}\) to contain terms with \emph{fewer} than two
derivatives of the metric. The freedom to use first derivatives does not
allow any new terms in \(\mathcal G_{\mu\nu}\) (see
Section~\ref{sec:6.1}), but if we can use the metric tensor itself, then one
new term is possible, equal to \(g_{\mu\nu}\) times a constant \(\Lambda\).
The field equations would then read
\[
  R_{\mu\nu}
  -\frac12g_{\mu\nu}R
  +\Lambda g_{\mu\nu}
  =
  8\pi G T_{\mu\nu}.
\]
The term \(\Lambda g_{\mu\nu}\) was originally introduced by
Einstein~\hyperref[ch7-ref-1]{[1]} for cosmological reasons (which have since
disappeared); for this reason, \(\Lambda\) is called the
\emph{cosmological constant}. This term satisfies the requirements (A), (C),
and (D), but does not satisfy (E), so \(\Lambda\) must be very small so as
not to interfere with the successes of Newton's theory of gravitation.
Except in Chapter~16, I am assuming throughout this book that
\(\Lambda=0\).
